\frontchapter{Préface}
\origpage{5}{前言 001}
国内外学界早已公认，语言学是以人类语言为研究对象的一门科学。从18世纪末历史比较语言学兴起，语言学大致经过了索绪尔的结构主义语言学、乔姆斯基的形式主义语言学和当代的功能主义语言学三个阶段，上下200多年历史，现已成为“一门领先的学科”\origfoot{1}{杨亦鸣，徐杰．语言学应该调整为一级学科 [J]．语言科学，2010（01）：5–13。}。这正说明语言学同其他任何科学一样，都是历史的产物。虽然并非每个语言工作者，特别是使用者都必须研究语言学史，但作为语言学者，包括专业的语言学习者还是应该在掌握学科理论的同时对其历史沿革有所了解，有所学习。正是基于这样的共识，我们注意到，国内许多高校的外语院系，包括法语院系，在本科高年级阶段就开设了语言学及其发展史的课程。而近年来，国内更有一些高校已经成立或正在酝酿成立语言学系，旨在培养具备全面语言学素养的专门人才，以适应今后语言学发展的需要。

与此同时，我们也毋庸讳言，语言学教学大多处于一个比较尴尬的境地，首先是因为人们心中的刻板印象，始终觉得语言学“艰深”、“晦涩”、“抽象”、“难懂”，而且“实用价值不大”，所以对它不是避而远之，便是敬而远之。其次是因为缺少真正对路的，契合我国广大青年学生心智特征、认知水平和求知需求的入门性普及型语言学教材，此类法语教科书目前在国内更是凤毛麟角。不过，同样不可否认的是，国内外越来越多的语言学家、学者和教师一直在为促使语言学这门“高尖深”的学问走下“神坛”，走出“象牙塔”而竭尽全力，并且取得了相当可喜的成就\origfoot{2}{恕我孤陋寡闻，挂一漏万，这里国内首先想到蓝纯的《语言学概论》（2009，外语教学与研究出版社）、姚小平的《西方语言学史》（2011，外语教学与研究出版社）和刘润清的《西方语言学流派》（1995，外语教学与研究出版社）。}。这当中，现在北京外国语大学任教的Lionel Spinosa先生编著的这本《语言学导论：从柏拉图到乔姆斯基》应当是最新的重要成果。

本书的第一个新在于，这是国内出版发行的第一部由法国教师根据中国法语专业学习者的兴趣、需求和能力用法语编写的普通语言学入门书，是我国高校法语专业学生了解并熟悉语言学的法语基本术语和概念的必备而重要的参照，为他们今后的专业深造与提升奠定良好的学术阅读基础。

本书的第二个新在于，语言风格上总体通俗易懂，娓娓道来，没有了一般语言学史著者的滞重之感，读来晓畅通达，生动明晰。这方面尤以其导读（Au lecteur）和开篇两章（Chapitres liminaires）的文字引人入胜，一下子拉近了作者与读者的距离，拉近了读者与作品的距离。
\origpage{6}{前言 002}
本书的第三个新在于，结构上思路清晰，全书以时间为轴线，以主题为线索，分三部分（语言学作为新兴科学的源起、语言学学科的6个领域，以及语言学学科的最新发展）介绍、阐述、解读了西方语言学自古希腊时期产生直至20世纪中叶的发展脉络，纵横2000年之久。诚如作者所言，阅读本书将是一次“穿越时空之旅”，既可了解语言学这一学科在各个发展时期的主要学术观点、标志性理论成果和主流思想流派，又能见识语言学的历史天空中灿若明星的语法学家、语言学家、逻辑学家和哲学家，其中有古印度人、古希腊人和古罗马人，也有现代的英国人、美国人、德国人、法国人、丹麦人、瑞士人，甚至中国人。内容不失丰足，而叙述洗练，同时在每一章的末尾辅以形式丰富的读后练习，帮助温故知新，激发掩卷思考，驱动深入研读。凡此种种，相信一定为广大学习和研究语言的专业人士以及对这一学科感兴趣的普通法语读者所接受并认可。这样的编排和设计非常便于学习者的自主阅读与学习，您可以时间为序，从头至尾通读全书，系统学习和了解西方语言学的历史沿革；您也可按主题索骥，直接查证和研读某一章节、某一流派或者某一代表人物。总之，开卷有益。

最后，我想特别提及的是，本书作者Lionel Spinosa副教授是热爱教书、关注学生、潜心教学的有心人。这部作品是他数年课堂教学、观察、反思和改进之后的产物，学生是他写作的不竭动力，完全可以说，这是一部献给中国高校法语专业各层级学生的爱心之作。

时值《语言学导论：从柏拉图到乔姆斯基》付梓，我热烈祝贺此书的出版，衷心祝贺我的法国同事和朋友作者Lionel的处女作终于问世，实可谓“十月怀胎，一朝分娩”。
\begin{flushright}傅\quad 荣\\2019/09/17 北京\end{flushright}
